% !TEX root = ../main.tex
\subsection*{8.2 分阶段实施路线与可验收成果}\label{page:40}
\noindent\textbf{实施路线以可复核产物和研究问题推进，而不以缺少依据的承诺填充时间表。}每一阶段均对应一个可以明确回答的问题：研究材料是否一致、数据能否进入同一主链、产业传导与证据是否可被评测、研究者能否在受控任务下使用系统、系统是否具备进入更大范围试用的治理条件。阶段之间的转换由验收材料触发，使后续数据和模型的扩展不会破坏既有报告的可追溯性。

\begin{center}\small 表40-1\quad 分阶段路线、交付物与验收门槛\par\vspace{1mm}
\begin{tabularx}{\linewidth}{p{0.16\linewidth}p{0.31\linewidth}X}
\toprule
阶段 & 核心交付物 & 阶段验收门槛\\
\midrule
阶段A：比赛材料固化 & 完整计划书、代码与数据版本索引、演示脚本、证据账本和逐页自检 & 文档可编译；图表、数据和引用可定位；演示固定到明确提交和\texttt{mock-v1}。\\
阶段B：数据主链整合 & 已登记的图谱、证据和评分版本，标准事件的完整报告 & 同一事件的路径、公司、评分和证据ID可以相互解析；异常输入形成可读失败记录。\\
阶段C：研究评测 & 固定事件任务集、人工标注、对照输出、消融记录和结果表 & 依页26协议报告样本、基线、指标、条件、分歧及原始记录。\\
阶段D：受控研究试用 & 明确角色和任务的试用脚本、操作记录与问题单 & 反馈可对应到版本、任务和具体界面，问题可被复现与分类处理。\\
阶段E：产品化准备 & 数据更新、权限、部署、监控和发布复核清单 & 完成数据授权、访问控制、性能、人工复核与运行治理的专项验收。\\
\bottomrule
\end{tabularx}\end{center}

阶段A保证比赛所展示的代码、图表、截图和叙事来自同一证据版本；阶段B把独立资产转化为统一报告能力；阶段C把报告能力置于可比较的研究任务中；阶段D检验研究者在实际流程中的使用体验；阶段E才讨论更广范围运行所需要的数据与工程治理。这条路线将系统成长拆分为连续的研究证据阶段，使每一步都能够用清晰的输入、输出和验收条件来定义。

项目负责人维护版本门禁和问题单：新增数据、规则、模型输出或界面能力在进入演示主线前，均需登记来源、接口、责任人、验证条件和影响范围。该机制一方面保留快速迭代空间，另一方面确保比赛展示中的任意页面都可回到对应提交、数据版本和复核记录。
